\chapter*{Agradecimientos}
\addcontentsline{toc}{chapter}{Agradecimientos}

A mi mamá, Verónica, y a mi papá, Américo, por el esfuerzo enorme que hicieron para pagarme parte de la carrera. Sé todo lo que se sacrificaron para darnos siempre lo mejor, y lo valoro profundamente. Soy la segunda ingeniera de la familia de mi papá y la primera de la de mi mamá, y espero que hoy se sientan tan orgullosos como yo de haber llegado hasta acá. Les agradezco por haber hecho siempre, desde el amor, lo que creyeron mejor para nosotros. Gracias por escuchar mis catarsis de la carrera y por decirme siempre las palabras justas para prepararme para el cuatrimestre siguiente.

A mi hermano, Nicolás, por inspirarme a estudiar esta carrera. Fue él quien me animó a cambiar de rumbo y empezar con Informática, me acompañó a inscribirme e hicimos juntos el curso de ingreso. Sin ese empujón, nada de esto habría empezado.

A mi novio, Braulio (Pipi), por acompañarme en cada etapa. Me explicó conceptos de la carrera, se sentó conmigo a ayudarme a entenderlos, me esperó con sushi cuando volvía cansada o de rendir, me cocinó en épocas de parciales y me tuvo paciencia cuando andaba estresada. Gracias por ser tan compañero y por tanto amor: tu apoyo incondicional fue clave para recorrer los últimos años de la carrera.

A toda mi familia, a Verónica, Américo, Nicolás, mi hermana Celeste y Braulio, por comprenderme y apoyarme, y por entender las veces que no pude estar en las juntadas familiares.

A mi gran amigo y compañero de tesis, Tomás Goncalves. Desde que nos conocimos en 2022 decidimos hacer todas las materias juntos, y fue mi soporte emocional durante toda la carrera. Hacerla con un compañero la hace mucho más llevadera, y compartir tanto tiempo con él fue un placer: pasó de ser mi amigo de la facu a ser mi mejor amigo. Estoy orgullosa de haber recorrido este camino a su lado. Lo quiero mucho y le agradezco por acompañarme en este proyecto y en toda la carrera.

A los profesores que me transmitieron sus conocimientos, que me inspiraron y que me hicieron querer tanto esta carrera.

A la Ing. Macarena Pomerantz, un agradecimiento especial. Me recomendaron tanto cursar Seminario 1 con ella que terminé pidiendo permiso en el trabajo para hacerlo por la mañana, y fue un placer aprender tanto a su lado. Además, me dio el gusto de ser la tutora de esta tesis, un proyecto al que le pusimos mucho amor. Maca trata a todos con tanta calidez que dan ganas de escucharla, aprender y hasta de que te corrija. Es una persona muy comprometida, y que fuera nuestra tutora me dio mucha seguridad, porque sabía que iba a estar para responder a tiempo y aportar su mirada, que es muy importante para mí.

A los amigos que hice durante la carrera, Diego, Flor, Gabriel, Bruno, Juan Ignacio, Carla y Federico, por motivarme siempre y por ayudarme a sacar adelante las materias. Haber compartido este camino con ustedes lo hizo mucho más llevadero.

A mis amigos de toda la vida, por bancarme y comprender que durante estos años no pude verlos tanto como hubiera querido.

\begin{flushright}\textit{Romina Micaela Palomino Gayoso}\end{flushright}

\bigskip

A mi papá, Fabián, quien me hizo hincha de Boca, por el enorme soporte que me brindó durante todos estos años de cursada, haciendo el esfuerzo de pagarme la carrera y dándome todo lo que necesité para poder estudiar tranquilo. Gracias también por haber sido quien me inspiró a seguir una carrera universitaria, ya que aunque seguimos ramas muy distintas, fue la primera persona que me inculcó que la educación es el arma más poderosa.

A mi mamá, Mónica, porque siempre estuvo atenta a cómo me sentía y en los detalles del día a día: desde dejarme la comida lista cuando llegaba agotado de la facultad, hasta asegurarse de que nunca me faltara nada. A ambos les agradezco profundamente por haberme inculcado el valor de la educación y por enseñarme a no bajar nunca los brazos. Gracias por creer en mí y espero que se sientan igual de orgullosos de mí de lo que yo estoy orgulloso de que ustedes sean mi padre y madre.

A mi hermana, Abril, a quien quiero muchísimo, y que aunque tenemos formas de pensar muy distintas, fue una gran fuente de inspiración para mí. No solo la admiro como la gran psicóloga que es, sino que verla en su recibida fue mi motivación y me inspiró a no rendirme en la recta final de la carrera, justo cuando ya sentía el cansancio de tantos años de estudio.

A mis mascotas, que las amo a todas: a Calu y Ollie, y a Susie y Luna, que aunque ya partieron, las amé y las amaré siempre con todo mi ser. Gracias por haber sido mi compañía silenciosa e incondicional en esas mañanas, tardes y noches en las que estudiaba solo frente a mi escritorio.

A mi compañera de tesis y mejor amiga, Romina Micaela Palomino Gayoso. Si hoy estoy en este lugar, es en gran parte gracias a ella. Mica fue y es un verdadero punto de referencia para mí; no miento si digo que en varios momentos pensé que, el día de mañana, quiero ser igual de dedicado, inteligente y profesional que ella. Definitivamente no habría llegado hasta acá sin su apoyo constante.

A los profesores que tuve a lo largo de toda mi carrera académica, quienes me transmitieron sus conocimientos y me dieron las herramientas necesarias para convertirme en el gran ingeniero que aspiro a ser.

A la Ing. Macarena Pomerantz, nuestra tutora de tesis. Aunque lamentablemente no tuve la suerte de cursar una materia con ella, pese a las excelentes referencias que siempre escuché de mis compañeros, estoy muy agradecido por toda la ayuda, la guía y la predisposición que nos brindó durante todo el año que llevó este proyecto. No me imagino una mejor tutora para esta etapa.

A las grandes personas que conocí durante la cursada y que hoy tengo la suerte de llamar amigos: Juan Ignacio, Diego, Bruno, Gabriel, Flor y Federico, por compartir este camino y hacer que todo fuera mucho más llevadero.

A mis amigos de toda la vida, esos que conozco desde la primaria y la secundaria, por haberme dado un apoyo enorme en cada etapa y por haberme sacado siempre una sonrisa en los momentos difíciles.

\begin{flushright}\textit{Tomás Gonçalves}\end{flushright}
